\chapter{\IfLanguageName{dutch}{Stand van zaken}{State of the art}}%
\label{ch:stand-van-zaken}

Dit hoofdstuk situeert de Gentse Feesten als casus voor lokale druktemeting. Eerst worden de veiligheidscontext en de beperkingen van bestaande meetmethoden besproken. Daarna volgt een afbakening van computervisiemethoden, met bijzondere aandacht voor \emph{density estimation}. De bespreking leidt tot de keuze van CSRNet als reproduceerbare eerste baseline.

\section{De Gentse Feesten als casus}

De Gentse Feesten ontstonden in 1843 als centraal georganiseerde Gemeentefeesten. Na een periode van teruglopende belangstelling kregen zij vanaf 1970 een nieuwe impuls rond Sint-Jacobs, onder meer door Walter De Buck en Trefpunt. Sindsdien groeiden de Feesten uit tot een groot cultureel volksfeest dat gedurende tien dagen in de Gentse binnenstad plaatsvindt. In 2021 werden de Gentse Feesten erkend als Vlaams Immaterieel Erfgoed. Tijdens de editie van 2023 werden meer dan 1,6 miljoen bezoekers geteld \autocite{StadGentGeschiedenis}.

In tegenstelling tot een evenement op een afgesloten terrein zijn de Gentse Feesten verspreid over pleinen, straten en binnenlocaties. Bezoekers kunnen zich vrij tussen die locaties verplaatsen en de bezetting verandert voortdurend doorheen de dag en nacht. Het totale bezoekersaantal zegt daarom weinig over de drukte op één specifieke plaats. Wanneer programma's gelijktijdig eindigen of bezoekersstromen samenkomen in smalle straten, kan de lokale bezetting snel toenemen.

Recente edities illustreren dat zulke veranderingen operationele aandacht vragen. Tijdens de openingsavond van 2023 liet de politie tijdelijk geen bijkomende bezoekers toe op Boomtown en het Laurentplein, omdat het er te druk werd \autocite{Schouppe2023}. In 2024 bleven de pleinen uitzonderlijk een uur langer open om een gelijktijdige uitstroom na een bijzonder drukke dag te vermijden \autocite{Aneca2024}. Deze situaties tonen niet noodzakelijk een incident aan, maar wel dat organisatoren soms snel moeten beslissen op basis van veranderende bezoekersstromen.

Stad Gent onderzocht eerder al, samen met onder meer de Universiteit Gent, verschillende vormen van druktemeting. Het project VLOED noemt fijnmazige informatie over de ruimtelijke spreiding van drukte, aangevuld met actuele gegevens en voorspellingen, als belangrijke behoefte \autocite{StadGentVloed}. Een automatische analyse van camerabeelden kan daarop aansluiten, op voorwaarde dat de resultaten begrijpelijk, betrouwbaar en zorgvuldig ingebed zijn in de bestaande besluitvorming.

\section{Risico's bij lokale drukte}

De veiligheid van een menigte hangt niet alleen af van het aantal aanwezigen. Ook de beschikbare ruimte, de vorm van doorgangen, de richting en snelheid van stromen, de duur van de drukte en de communicatie tussen betrokken diensten spelen een rol \autocite{Helbing2000,Zeitz2009,Hutton2025}. Hoge dichtheden kunnen de bewegingsvrijheid sterk beperken; in combinatie met ongunstige omstandigheden kunnen zij tot gevaarlijke druk en verlies van controle over de beweging leiden. Een systeem voor druktemeting moet daarom als vroegtijdige informatiebron worden beschouwd, niet als een zelfstandige voorspeller van incidenten.

Literatuur noemt dichtheden van ongeveer acht tot tien personen per vierkante meter als bijzonder kritisch \autocite{Haghani2022}. Een vaste alarmwaarde kan echter niet rechtstreeks uit een camerabeel of density map worden afgeleid. Daarvoor moet elk beeld eerst via camerakalibratie en een grondvlaktransformatie aan een reële oppervlakte worden gekoppeld. Ook de locatie, de richting van de stroom en de duur van een verhoogde dichtheid moeten in een operationele drempel worden meegenomen. In deze bachelorproef wordt daarom geen veiligheidsdrempel ingesteld; de PoC beperkt zich tot personentellingen en relatieve druktepatronen in beeldcoördinaten.

\section{Bestaande meetmethoden en de rol van automatisering}

Organisatoren combineren vaak terreinobservaties, handmatige tellingen en ervaring uit voorgaande edities. Deze bronnen blijven waardevol, maar zijn niet continu en niet overal tegelijk beschikbaar. Tussen de waarneming, de interpretatie, de communicatie en een eventuele maatregel ontstaat bovendien onvermijdelijk vertraging. Automatische metingen kunnen die menselijke beoordeling aanvullen met een regelmatige en reproduceerbare indicatie van lokale veranderingen.

WiFi- en Bluetoothsignalen kunnen worden gebruikt om personen te tellen of te lokaliseren \autocite{GuillenPerez2019}. Hun ruimtelijke resolutie en nauwkeurigheid hangen echter af van de meetopstelling, de aanwezigheid van apparaten en de gebruikte methode. Camerabeelden bieden in principe een fijnere ruimtelijke observatie van één afgebakende zone, maar stellen eisen aan zichtbaarheid, rekenkracht, privacy en robuustheid tegen variaties in de scène. Computervisie is daarom geen vervanging van andere meetbronnen, maar een mogelijke aanvullende sensor.

\section{Computervisie voor crowd counting}

De meeste moderne methoden voor beeldanalyse gebruiken convolutionele neurale netwerken (CNN's), die visuele patronen uit beelden leren \autocite{OShea2015}. Binnen crowd counting zijn twee uitvoerwijzen relevant. \emph{Objectdetectie} zoekt individuele personen en voorspelt voor elke detectie een begrenzingskader. \emph{Density estimation} voorspelt daarentegen een tweedimensionale kaart waarvan de waarden de lokale aanwezigheid van personen weergeven. Beide benaderingen kunnen CNN's gebruiken; het verschil ligt dus in de taak en de uitvoer, niet in het al dan niet gebruiken van een CNN.

Objectdetectie is aantrekkelijk wanneer personen voldoende zichtbaar en van elkaar te onderscheiden zijn. YOLO is bijvoorbeeld ontworpen voor snelle objectdetectie \autocite{Redmon2016}. In dichte menigtes verminderen overlap, schaalverschillen en gedeeltelijke zichtbaarheid echter de betrouwbaarheid van afzonderlijke detecties \autocite{Khan2020}. Een density-estimationmodel hoeft niet voor elke persoon een afzonderlijk kader te vinden. Het kan daardoor in zulke scènes een bruikbare telling en ruimtelijk patroon leveren, al blijft de nauwkeurigheid afhankelijk van de trainingsdata en de gelijkenis met de doelomgeving.

Een \emph{density map} is geen directe kaart in personen per vierkante meter. De kaart is gekoppeld aan beeldpixels of aan cellen in de uitvoerresolutie van het model. De som van alle waarden levert een geschat aantal personen in het zichtveld. Pas na camerakalibratie en projectie naar het grondvlak kan een dichtheidswaarde aan een fysieke oppervlakte worden gekoppeld.

\section{Selectie van een eerste baseline}

Tabel~\ref{tab:long-list-modellen} vergelijkt kandidaatmodellen op basis van hun hoofdidee, hun sterkte voor deze casus en hun voornaamste beperking. De tabel is geen benchmark: de modellen worden niet onder identieke omstandigheden uitgevoerd. Zij dient uitsluitend om de keuze voor een eerste implementatie te verantwoorden.

\begin{table}[htbp]
    \centering
    \small
    \begin{tabular}{|p{2.4cm}|p{2.6cm}|p{4.2cm}|p{4.1cm}|}
        \hline
        \textbf{Model} & \textbf{Categorie} & \textbf{Relevante sterkte} & \textbf{Beperking voor deze PoC} \\
        \hline
        YOLO \autocite{Redmon2016} & Objectdetectie & Snelle detectie wanneer personen afzonderlijk zichtbaar zijn. & Overlap tussen personen bemoeilijkt afzonderlijke detecties. \\
        \hline
        Faster R-CNN & Objectdetectie & Sterke referentie voor nauwkeurige objectdetectie. & Complexer en doorgaans trager voor een eerste baseline met real-timeambitie. \\
        \hline
        SSD & Objectdetectie & Relatief compacte detector met een bruikbare snelheids-nauwkeurigheidsbalans. & Blijft afhankelijk van afzonderlijke begrenzingskaders. \\
        \hline
        RetinaNet & Objectdetectie & Ontworpen om de onbalans tussen achtergrond en objecten te beperken. & Niet specifiek gericht op zeer dichte menigtes. \\
        \hline
        CSRNet \autocite{Li2018} & Density estimation & Gedilateerde convoluties behouden context en ruimtelijk detail. & Vereist training en validatie voor de doelomgeving. \\
        \hline
        MCNN \autocite{Zhang2016} & Density estimation & Meerdere kolommen behandelen schaal- en perspectiefvariatie. & Oudere en zwaardere architectuur als eerste baseline. \\
        \hline
        SANet & Density estimation & Combineert informatie op verschillende schalen. & Hogere implementatiecomplexiteit. \\
        \hline
        CAN & Density estimation & Gebruikt contextuele informatie voor de dichtheidsschatting. & Hogere complexiteit en minder geschikt voor een minimale baseline. \\
        \hline
    \end{tabular}
    \caption[Selectie van kandidaatmodellen]{Kwalitatieve selectie van kandidaatmodellen voor druktemeting.}
    \label{tab:long-list-modellen}
\end{table}

CSRNet wordt geselecteerd als eerste baseline omdat de architectuur specifiek is ontwikkeld voor dicht bezette scènes, volledig convolutioneel is en goed gedocumenteerd werd op ShanghaiTech \autocite{Li2018}. MCNN blijft een zinvolle referentie voor toekomstig vergelijkend onderzoek, maar wordt binnen de beschikbare scope niet geïmplementeerd. De resultaten in deze bachelorproef mogen bijgevolg niet worden gelezen als een vergelijking tussen modellen.

\section{Privacy en gegevensbescherming}

Density estimation heeft als voordeel dat de PoC geen gezichten herkent, personen heridentificeert of trajecten volgt. Dat beperkt de functionaliteit tot een telling en een ruimtelijk patroon. De camerabeelden kunnen desondanks personen identificeerbaar in beeld brengen en zijn dan persoonsgegevens. Een density-estimationmodel maakt de verwerking dus niet automatisch anoniem of conform de AVG.

Een productie-oplossing vereist minstens een welbepaalde rechtsgrond en doel, dataminimalisatie, transparante informatie voor bezoekers, beperkte bewaartermijnen, toegangscontrole en passende beveiliging. Bij grootschalige monitoring van een publiek toegankelijke ruimte kan bovendien een gegevensbeschermingseffectbeoordeling vereist zijn \autocite{EDPBVideo2020}. Praktische maatregelen, zoals verwerking aan de rand van het netwerk, het niet bewaren van ruwe beelden en het beperken van de analyse tot afgebakende zones, verdienen daarbij verder onderzoek. De PoC evalueert deze maatregelen niet en doet geen uitspraak over juridische conformiteit.
